\documentclass[journal]{IEEEtran}

% ============================================================
% PACKAGES
% ============================================================
\usepackage{amsmath,amssymb}
\usepackage{booktabs}
\usepackage{array}
\usepackage{multirow}
\usepackage{cite}
\usepackage{url}
\usepackage{microtype}
\usepackage{xcolor}

% ============================================================
% TABLE COLUMN TYPES
% ============================================================
\newcolumntype{L}[1]{>{\raggedright\arraybackslash}p{#1}}
\newcolumntype{C}[1]{>{\centering\arraybackslash}p{#1}}

% ============================================================
% DOCUMENT
% ============================================================
\begin{document}

% ============================================================
% TITLE
% ============================================================

\title{Vision-Guided Multi-Objective Reinforcement Learning\\
for UR5 Manipulation in Dynamic Environments}

\author{Research~Draft}

\maketitle

\begin{center}
\small
\textit{Sections 2--4: Related Work, Contributions, and Expected Results}
\end{center}

\vspace{0.5em}

% ============================================================
% SECTION 2
% ============================================================

\section{Related Work}

Real-world robot manipulation requires a robot to achieve its task
while simultaneously considering conflicting objectives such as task
success, energy, smoothness, and safety. Existing research has made
strong progress in multi-objective reinforcement learning (MORL) and
vision-guided manipulation, but these areas have largely developed
separately.

The proposed work focuses on the combination of five capabilities:

\begin{enumerate}
    \item \textbf{Vision:} RGB-D / point-cloud observations.
    \item \textbf{MORL:} vector rewards and a Pareto-optimal solution set.
    \item \textbf{Manipulation:} contact-rich 6-DoF arm control.
    \item \textbf{Preference conditioning:} runtime modification of
    behavior using a preference vector $\lambda$ without retraining.
    \item \textbf{Safety as a constraint:} collision, velocity, and
    workspace limits formulated as a constrained MDP (CMDP).
\end{enumerate}

Table~\ref{tab:capability} provides the high-level comparison.

% ============================================================
% TABLE 1
% ============================================================

\begin{table*}[t]
\centering
\caption{Capability comparison between representative related work and
the proposed approach.}
\label{tab:capability}

\renewcommand{\arraystretch}{1.25}
\setlength{\tabcolsep}{4pt}

\begin{tabular}{l c c c c c c}

\toprule
\textbf{Method} &
\textbf{Vision} &
\textbf{MORL} &
\textbf{Dynamic} &
\textbf{Manip.} &
\textbf{Preference} &
\textbf{Safety Constraint} \\
\midrule

Envelope MOQ (NeurIPS 2019)
& $\times$ & \checkmark & $\times$ & $\times$ & \checkmark & $\times$ \\

PGMORL (ICML 2020)
& $\times$ & \checkmark & $\times$ & $\times$ & \checkmark & $\times$ \\

MO-MPO (NeurIPS 2022)
& $\times$ & \checkmark & $\times$ & partial & \checkmark & $\times$ \\

DOL / ImageDeepSea
& partial & \checkmark & $\times$ & $\times$ & \checkmark & $\times$ \\

Diffusion Policy
(RSS 2023 / IJRR 2024)
& \checkmark & $\times$ & \checkmark & \checkmark & $\times$ & partial \\

DP3 (RSS 2024)
& \checkmark & $\times$ & \checkmark & \checkmark & $\times$ & empirical \\

ACT / ALOHA
& \checkmark & $\times$ & \checkmark & \checkmark & $\times$ & $\times$ \\

GCR-PPO (2025)
& $\times$ & \checkmark & partial & \checkmark & $\times$ & $\times$ \\

PRC-CMORL
(Neurocomputing 2023)
& $\times$ & \checkmark & $\times$ & $\times$ & \checkmark & \checkmark \\

Demo-enhanced MORL navigation
(IROS 2025)
& lidar & \checkmark & \checkmark & $\times$ & \checkmark & $\times$ \\

POVNav (2025)
& \checkmark & classical & \checkmark & $\times$ & partial & \checkmark \\

\textbf{This work}
& \textbf{\checkmark} & \textbf{\checkmark} &
\textbf{\checkmark} & \textbf{\checkmark} &
\textbf{\checkmark} & \textbf{\checkmark} \\

\bottomrule
\end{tabular}
\end{table*}

% ============================================================
% 2.1
% ============================================================

\subsection{Preference-Conditioned MORL}

The MORL literature provides the foundation for learning policies that
balance multiple objectives.

\textbf{Envelope MOQ-learning} \cite{yang2019} trains a single network
$Q(s,a,\omega)$ over a preference space. Its envelope Bellman operator
and hindsight experience replay provide the foundation for
single-policy preference conditioning.

\textbf{PGMORL} \cite{xu2020} instead maintains a population of PPO
policies trained with different preference vectors. A Gaussian-process
prediction model guides the population toward the Pareto front.
Hypervolume and sparsity are used to evaluate the resulting solutions.

\textbf{MO-MPO} \cite{abdolmaleki2022} uses per-objective KL budgets
$\varepsilon_k$ instead of direct scalar weighting, providing a
scale-invariant treatment of multiple objectives.

\textbf{DOL} \cite{mossalam} demonstrates multi-objective learning with
image observations through the ImageDeepSea environment. However, the
visual environment is a toy grid rather than a camera-based robotic
manipulation problem.

\textbf{MO-Gymnasium} \cite{felten2023} provides a standard vector-reward
interface for multi-objective environments.

\textbf{Main limitation:} these methods provide the MORL foundation,
but do not simultaneously address visual 6-DoF manipulation in dynamic
scenes with explicit safety constraints.

% ============================================================
% 2.2
% ============================================================

\subsection{Vision-Guided Robot Manipulation}

Vision-guided manipulation has developed rapidly through imitation
learning and visuomotor policy learning.

\textbf{Diffusion Policy} \cite{chi2023} represents the policy as a
conditional diffusion process over action sequences. Its reported UR5
setup uses $T_o=2$, $T_p=16$, and $T_a=8$, with a 10 Hz policy and
125 Hz interpolation. The reported velocity limit is $<0.43$ m/s and
the workspace maintains a position $\geq 1$ cm above the table.

\textbf{3D Diffusion Policy (DP3)} \cite{ze2024} uses sparse point-cloud
representations instead of RGB images. The method uses 512 or 1024
points and reports a 24.2\% relative gain over image baselines on
72 simulation tasks, together with 85\% real success with 40 demos.

\textbf{ACT / ALOHA} \cite{zhao2023},
\textbf{Mobile ALOHA} \cite{fu2024},
\textbf{UMI} \cite{chi2024}, and
\textbf{IRIS} \cite{jiang2025} further demonstrate the effectiveness of
demonstration-based visuomotor manipulation.

These methods provide the vision and manipulation foundation for the
proposed system.

\textbf{Main limitation:} these approaches primarily optimize a
single task objective. They do not provide a runtime preference
mechanism for changing the trade-off among multiple objectives.

% ============================================================
% 2.3
% ============================================================

\subsection{Robot-Scale MORL}

\textbf{GCR-PPO} \cite{munn2025} provides an important algorithmic
foundation for the proposed multi-objective critic.

The method uses a multi-head critic,

\begin{equation}
V_{\phi}:S\rightarrow\mathbb{R}^{K},
\end{equation}

with component-wise GAE, ratio-preserving advantage normalization, and
asymmetric PCGrad.

GCR-PPO reports $\sim$4.55\% average return gain over PPO, with up to
9.5\% on custom multi-objective suites.

These reported values are used only as an \textbf{algorithmic reference}
for the scale of improvement observed in the source study. They are
not adopted as numerical return targets for the proposed UR5 system.

However, its observations are proprioceptive and the method does not
combine a visual encoder, runtime preference conditioning, or explicit
CMDP safety.

% ============================================================
% 2.4
% ============================================================

\subsection{Safety-Constrained MORL}

\textbf{PRC-CMORL} \cite{he2023} formulates multi-objective robot
control as a constrained multi-objective MDP. It provides a foundation
for treating safety as a constraint rather than simply another
objective.

This distinction is important:

\begin{center}
\textbf{Safety objective} $\neq$ \textbf{Safety constraint}.
\end{center}

If safety is included only as a reward term, the policy may trade safety
for task performance when $\lambda$ changes. In contrast, a CMDP
explicitly requires safety constraints to remain satisfied within the
constraint formulation.

The proposed work therefore uses safety constraints such as

\begin{equation}
d(s,O)\geq D_{\mathrm{safe}},
\end{equation}

together with

\begin{equation}
\|\dot{x}\|<0.43~\mathrm{m/s}.
\end{equation}

Safety is therefore handled separately from the preference vector.
The preference vector controls the task-objective trade-off among
success, energy, and smoothness.

% ============================================================
% 2.5
% ============================================================

\subsection{Vision and MORL: The Closest Existing Work}

\textbf{Demonstration-enhanced adaptable MO robot navigation}
\cite{deheuvel2025} is the closest intersection of MORL,
preference conditioning, dynamic environments, and sim-to-real.

It uses a preference interpolator

\begin{equation}
I(\lambda)=\lambda^p
\end{equation}

and adapts the preference at runtime without retraining.

However, the system uses lidar and wheeled robot navigation rather than
visual 6-DoF manipulation. Collision is also treated as a sparse
objective rather than a CMDP constraint.

\textbf{POVNav} \cite{pushp2025} uses Pareto reasoning together with
visual information and explicit safety constraints, but its focus is
navigation and classical planning rather than a learned
preference-conditioned manipulation policy.

% ============================================================
% 2.6
% ============================================================

\subsection{Research Gap}

The literature therefore provides individual building blocks:

\begin{table*}[t]
\centering
\caption{What existing research provides and what remains to be
combined.}
\label{tab:gap}

\renewcommand{\arraystretch}{1.2}

\begin{tabular}{l l l}
\toprule
\textbf{Research Area}
& \textbf{What It Provides}
& \textbf{Remaining Gap} \\
\midrule

MORL
& Pareto optimization
& Limited visual manipulation \\

Vision manipulation
& Visual 6-DoF control
& Mostly single-objective \\

GCR-PPO
& Multi-objective critic
& No visual preference-conditioned UR5 \\

PRC-CMORL
& Safety constraints
& No visual manipulation \\

IROS 2025 MORL navigation
& Preference adaptation + dynamic scenes
& Navigation rather than manipulation \\

POVNav
& Visual Pareto planning + safety
& Classical navigation planning \\

\textbf{This work}
& \textbf{Combines all five}
& \textbf{UR5 dynamic manipulation} \\

\bottomrule
\end{tabular}
\end{table*}

The resulting research gap is the combination of:

\begin{center}
\textbf{Vision + MORL + 6-DoF Manipulation +
Preference Conditioning + CMDP Safety}.
\end{center}

The proposed work investigates this combination on a UR5 in static and
dynamic environments.

% ============================================================
% SECTION 3
% ============================================================

\section{Contributions}

The proposed study has three main contributions.

% ------------------------------------------------------------
% C1
% ------------------------------------------------------------

\subsection*{C1. Preference-Conditioned Vision-Guided MORL}

The first contribution is a vision-guided MORL framework for UR5
manipulation in which the policy can change its behavior according to a
runtime preference vector $\lambda$ without retraining.

The overall architecture is

\begin{equation}
\boxed{
\begin{aligned}
\text{RGB-D}
&\rightarrow
\text{Point Cloud}
\rightarrow
\text{DP3}\\
&\rightarrow
\text{MO Critic}
\rightarrow
\text{Diffusion Policy}
\rightarrow
\text{UR5}
\end{aligned}
}
\end{equation}

while safety is handled separately through CMDP constraints.

The multi-objective critic uses component-wise GAE,
ratio-preserving normalization, and asymmetric PCGrad.

% ------------------------------------------------------------
% C2
% ------------------------------------------------------------

\subsection*{C2. Multi-Objective UR5 Benchmark}

The second contribution is a vision-based multi-objective UR5
benchmark containing both static and dynamic environments.

The benchmark considers:

\begin{table}[t]
\centering
\caption{Proposed UR5 multi-objective benchmark structure.}
\label{tab:benchmark}

\renewcommand{\arraystretch}{1.2}

\begin{tabular}{l l}
\toprule
\textbf{Category} & \textbf{Proposed Setting} \\
\midrule

Robot & UR5 \\

Observation & RGB-D / point cloud \\

Objectives & Success, energy, smoothness/jerk \\

Safety & CMDP constraint \\

Static scene & Static clutter \\

Dynamic scene 1 & Moving obstacle \\

Dynamic scene 2 & Moving target \\

Dynamic scene 3 & Sudden appearance \\

Preference & Runtime $\lambda$ \\

\bottomrule
\end{tabular}
\end{table}

The benchmark is wrapped using the MO-Gymnasium API so that existing
MORL baselines can be evaluated using the same environment definition.

% ------------------------------------------------------------
% C3
% ------------------------------------------------------------

\subsection*{C3. Controlled Evaluation and Sim-to-Real Study}

The third contribution is an evaluation protocol designed to separate
the effects of:

\begin{enumerate}
    \item multi-objective learning,
    \item preference conditioning,
    \item visual representation, and
    \item safety constraints.
\end{enumerate}

The main comparisons are:

\begin{table}[t]
\centering
\caption{Main experimental comparisons.}
\label{tab:comparisons}

\renewcommand{\arraystretch}{1.2}

\footnotesize

\begin{tabular}{l l l}
\toprule
\textbf{Comparison}
& \textbf{Baseline}
& \textbf{Purpose} \\
\midrule

MORL
& PGMORL
& Pareto performance \\

MORL
& Envelope MOQ
& Preference-conditioned MORL \\

Vision
& Single-objective DP3
& Visual manipulation reference \\

Objective
& Visual scalarized PPO
& Effect of vector reward \\

Preference
& Fixed $\lambda$
& Effect of runtime preference \\

Safety
& Collision-as-objective
& Effect of CMDP constraint \\

\bottomrule
\end{tabular}
\end{table}

Simulation and real-robot experiments are conducted in parallel so that
the sim-to-real gap can be explicitly measured.

% ============================================================
% SECTION 4
% ============================================================

\section{Expected Results and Comparison}

The proposed method is evaluated against published results and against
the internal baselines defined for this study. The numbers in the
``Expected Result'' column are \textbf{research targets or reference
values, not measured results}. Published results are used to motivate
the experimental design, while the proposed targets define the
hypotheses to be tested.

% ============================================================
% 4.1 OVERALL EXPECTED COMPARISON
% ============================================================

\subsection{Overall Expected Comparison}

Table~\ref{tab:overall_expected} provides the main comparison between
reported results from existing methods and the expected results of the
proposed Vision-Guided MORL framework.

\begin{table*}[t]
\centering
\caption{Overall expected comparison between published references and
the proposed method.}
\label{tab:overall_expected}
\renewcommand{\arraystretch}{1.35}
\setlength{\tabcolsep}{4pt}

\begin{tabular}{L{2.2cm} L{3.2cm} L{3.2cm} L{5.0cm}}
\toprule
\textbf{Method / Reference}
&
\textbf{Main Result Reported in the Literature}
&
\textbf{Our Method}
&
\textbf{Expected Result}
\\
\midrule

DP3 \cite{ze2024}
&
85\% real-robot success with 40 demonstrations
&
DP3 + MORL + preference conditioning
&
\textbf{$\geq70\%$ real static success target.}
The 85\% DP3 result is retained as a single-objective
external reference, not as the target of the complete MORL system.
\\

DP3 \cite{ze2024}
&
24.2\% relative gain of point-cloud representation over image
baselines
&
DP3 point-cloud encoder
&
Retain the point-cloud representation as the vision backbone;
the 24.2\% value is used as a representation reference, not as a
proposed-system improvement target.
\\

Diffusion Policy \cite{chi2023}
&
125 Hz UR5 control loop; velocity $<0.43$ m/s
&
Preference-conditioned diffusion policy
&
Maintain the 125 Hz control loop and the
$<0.43$ m/s Cartesian velocity constraint.
\\

GCR-PPO \cite{munn2025}
&
$\sim$4.55\% average return improvement; up to 9.5\% on selected
benchmarks
&
Multi-head MO critic + asymmetric PCGrad
&
Use the reported values only as an algorithmic
order-of-magnitude reference. The proposed UR5 system is not assigned
a 4.55\% or 9.5\% return target.
\\

PGMORL \cite{xu2020}
&
Multi-policy Pareto optimization
&
Single preference-conditioned policy
&
Hypervolume $\geq$ PGMORL on the same UR5 benchmark wrapper.
\\

Envelope MOQ \cite{yang2019}
&
Single-model preference-conditioned MORL
&
Vision + preference-conditioned MORL
&
Hypervolume $\geq$ Envelope MOQ on the same benchmark wrapper while
additionally handling visual manipulation.
\\

PRC-CMORL \cite{he2023}
&
Constrained MORL evaluated using Pareto-quality metrics
&
CMDP-constrained UR5 MORL
&
Strictly lower safety-violation rate than the
collision-as-objective baseline, with static violations approaching
0\%.
\\

Proposed Method
&
Not yet measured
&
DP3 + MO critic + preference-conditioned diffusion + CMDP
&
$\geq70\%$ real static success target; improved Pareto quality;
measurable preference-dependent behavior; and reduced safety
violations.
\\

\bottomrule
\end{tabular}
\end{table*}

% ============================================================
% 4.2 NUMERICAL TARGETS FOR EACH RESEARCH QUESTION
% ============================================================

\subsection{Numerical Targets for Each Research Question}

% ------------------------------------------------------------
% RQ1
% ------------------------------------------------------------

\subsubsection{RQ1 --- Can the visual representation support
multi-objective UR5 control?}

The first experiment verifies whether the DP3 point-cloud
representation provides sufficient information for learning the
individual objectives before combining them into MORL.

\begin{table*}[t]
\centering
\caption{RQ1 visual representation targets.}
\label{tab:rq1}
\renewcommand{\arraystretch}{1.25}

\begin{tabular}{l l l}
\toprule
\textbf{Metric}
&
\textbf{Reference / Baseline}
&
\textbf{Expected Result}
\\
\midrule

Point-cloud size
&
DP3: 512 or 1024 points
&
\textbf{512 points}
\\

Demonstrations per task
&
DP3: 10--20
&
\textbf{10--20 demonstrations} for warm-start
\\

Simulation success
&
Proposed visual benchmark
&
\textbf{Approximately 92\%}
\\

Real static success
&
DP3: 85\% with 40 demonstrations
&
\textbf{Target: $\geq70\%$}
\\

Point-cloud representation
&
DP3: 24.2\% relative gain
&
\textbf{Use point-cloud encoder}
\\

Objectives
&
Proposed benchmark
&
\textbf{Success, energy, smoothness}
\\

Safety
&
Conventional reward formulation
&
\textbf{Safety treated as a CMDP constraint}
\\

\bottomrule
\end{tabular}
\end{table*}

The approximately 92\% simulation result is used as the simulation
reference for the proposed visual manipulation setup. The real-robot
feasibility target is $\geq70\%$ static success.

The 85\% real-robot result reported by DP3 is retained as a
\textbf{single-objective external reference}. It is not transferred
directly to the complete MORL system.

If the individual visual objectives cannot be learned reliably before
MORL is introduced, the resulting limitation is likely related to the
visual observation or reward formulation rather than the MORL
algorithm itself.

% ------------------------------------------------------------
% RQ2
% ------------------------------------------------------------

\subsubsection{RQ2 --- Does the multi-head MORL critic improve
performance?}

GCR-PPO provides the algorithmic reference for the proposed
multi-objective critic.

However, the reported 4.55\% average and 9.5\% maximum improvements are
\textbf{not copied as numerical UR5 return targets}. They are used only
to establish the order of magnitude of improvement reported for a
related multi-objective optimization method.

\begin{table*}[t]
\centering
\caption{RQ2 multi-objective critic comparison.}
\label{tab:rq2}
\renewcommand{\arraystretch}{1.25}
\setlength{\tabcolsep}{4pt}

\begin{tabular}{l l l l}
\toprule
\textbf{Method}
&
\textbf{Average Return}
&
\textbf{Maximum Gain}
&
\textbf{Role in Our Study}
\\
\midrule

Scalarized PPO
&
Baseline
&
0\%
&
\textbf{Visual baseline}
\\

GCR-PPO
&
$\sim$104.55\% of PPO
&
$\sim$109.5\%
&
Literature reference
\\

Proposed MO critic
&
Not predetermined
&
Not predetermined
&
\textbf{Evaluate Pareto quality}
\\

\bottomrule
\end{tabular}
\end{table*}

The scalarized PPO baseline will use the same visual observation
setting as the proposed system. Therefore, the comparison does not
treat scalarized PPO as a non-visual baseline.

The main expected improvement is therefore evaluated using Pareto-front
quality rather than assuming a particular percentage return gain.

The primary metrics are:

\begin{itemize}
    \item hypervolume,
    \item Pareto coverage,
    \item task success,
    \item energy,
    \item smoothness, and
    \item safety-violation rate.
\end{itemize}

The GCR-PPO values of approximately 4.55\% and 9.5\% are therefore
\textbf{reference values only}, not guaranteed results for the proposed
UR5 system.

% ------------------------------------------------------------
% RQ3
% ------------------------------------------------------------

\subsubsection{RQ3 --- Does preference conditioning produce different
behaviors without retraining?}

The main purpose of preference conditioning is to allow the same
trained policy to produce different behaviors by changing the
preference vector $\lambda$.

The preference vector is defined as

\begin{equation}
\lambda =
[
\lambda_{\mathrm{success}},
\lambda_{\mathrm{energy}},
\lambda_{\mathrm{smoothness}}
].
\end{equation}

\textbf{Safety is not included in $\lambda$.} It is handled separately
through the CMDP formulation.

\begin{table*}[t]
\centering
\caption{RQ3 preference-conditioning comparison.}
\label{tab:rq3}
\renewcommand{\arraystretch}{1.25}
\setlength{\tabcolsep}{4pt}

\footnotesize

\begin{tabular}{l l l l}
\toprule
\textbf{Experiment}
&
\textbf{Existing Approach}
&
\textbf{Proposed Approach}
&
\textbf{Expected Result}
\\
\midrule

Policy training
&
Separate scalarized policies
&
\textbf{One policy}
&
One trained policy
\\

Preference
&
Fixed scalar weights
&
\textbf{Continuous $\lambda$}
&
Runtime-adjustable $\lambda$
\\

Retraining after preference change
&
Required in separate-policy setup
&
\textbf{Not required}
&
No retraining
\\

Number of tested preferences
&
Separate policies
&
$\lambda_1,\ldots,\lambda_5$
&
\textbf{3--5 settings}
\\

Behavior
&
Fixed trade-off
&
Preference-dependent
&
\textbf{Measurably different behavior}
\\

Pareto quality
&
Baseline
&
Proposed
&
\textbf{HV $\geq$ PGMORL / Envelope target}
\\

Real-robot transfer
&
Individual policies
&
Preference-conditioned policy
&
\textbf{Trade-offs observable on real UR5}
\\

\bottomrule
\end{tabular}
\end{table*}

For example, the preference settings may emphasize different
combinations of the three task objectives:

\begin{table*}[t]
\centering
\caption{Illustrative preference settings.}
\label{tab:preference_settings}
\renewcommand{\arraystretch}{1.2}

\footnotesize

\begin{tabular}{l l l}
\toprule
\textbf{Setting}
&
\textbf{Main Priority}
&
\textbf{Expected Robot Behavior}
\\
\midrule

$\lambda_1$
&
Task success
&
Higher success emphasis, potentially higher energy
\\

$\lambda_2$
&
Energy
&
Lower energy, potentially longer execution
\\

$\lambda_3$
&
Smoothness
&
Smoother trajectories and motion
\\

$\lambda_4$
&
Balanced
&
Balanced success--energy--smoothness behavior
\\

$\lambda_5$
&
Custom mixed
&
User-selected combination of the three objectives
\\

\bottomrule
\end{tabular}
\end{table*}

The important result is not that one preference is universally better.
The expected result is that

\begin{equation}
\boxed{
\begin{aligned}
\lambda_i &\neq \lambda_j\\
&\Rightarrow\quad \text{measurably different robot behavior}
\end{aligned}
}
\end{equation}

without retraining the policy.

Safety is deliberately excluded from this preference table because it
is represented by the CMDP constraints rather than by a preference
weight.

% ------------------------------------------------------------
% RQ4
% ------------------------------------------------------------

\subsubsection{RQ4 --- Does CMDP safety reduce safety violations?}

Safety is treated as a constraint instead of simply adding collision
to the reward vector.

The proposed safety constraints are

\begin{equation}
d(s,O)\geq D_{\mathrm{safe}},
\end{equation}

\begin{equation}
\|\dot{x}\| < 0.43~\mathrm{m/s},
\end{equation}

together with workspace limits.

\begin{table*}[t]
\centering
\caption{RQ4 safety comparison.}
\label{tab:rq4}
\renewcommand{\arraystretch}{1.2}
\setlength{\tabcolsep}{4pt}

\begin{tabular}{l l l l}
\toprule
\textbf{Safety Metric}
&
\textbf{Collision-as-Objective}
&
\textbf{Proposed CMDP}
&
\textbf{Expected Result}
\\
\midrule

Collision handling
&
Reward penalty
&
Explicit constraint
&
Lower violation rate
\\

Minimum obstacle distance
&
Reward dependent
&
$d(s,O)\geq D_{\mathrm{safe}}$
&
Lower violation rate
\\

Cartesian velocity
&
Not necessarily constrained
&
$<0.43$ m/s
&
\textbf{Below 0.43 m/s}
\\

Workspace
&
Reward dependent
&
Explicit bounds
&
Reduced workspace violations
\\

Static safety violation rate
&
Baseline measurement
&
CMDP
&
\textbf{Approaches 0\%}
\\

Dynamic safety violation rate
&
Collision objective
&
CMDP
&
\textbf{Lower than collision-as-objective}
\\

Task success
&
Optimized independently
&
Safety + task objectives
&
A small reduction is acceptable if safety improves
\\

\bottomrule
\end{tabular}
\end{table*}

The primary expected result is therefore

\begin{equation}
\boxed{
\begin{aligned}
\text{CMDP violation rate}\\
&<\text{collision-as-objective violation rate}
\end{aligned}
}
\end{equation}

with the static validation set expected to approach

\begin{equation}
\boxed{
\text{Static safety violation rate}
\approx 0\%
}
\end{equation}

The formulation does \textbf{not} claim that every safety constraint
will always be satisfied. The experimental criterion is a
\textbf{strictly lower violation rate} than the collision-as-objective
baseline, with static violations approaching zero.

A small reduction in task success is acceptable if it results in a
substantial reduction in safety violations. The objective is therefore
a better constrained Pareto set rather than independently maximizing
every metric.

% ============================================================
% 4.3 MAIN QUANTITATIVE TARGET TABLE
% ============================================================

\subsection{Main Quantitative Target Table}

The following table summarizes the numerical targets and literature
reference values that will be used to determine whether the proposed
method achieves its intended improvements.

\begin{table*}[t]
\centering
\caption{Main quantitative targets and literature reference values.}
\label{tab:quantitative_targets}
\renewcommand{\arraystretch}{1.3}
\setlength{\tabcolsep}{4pt}

\begin{tabular}{l l l l}
\toprule
\textbf{Performance Area}
&
\textbf{Published Reference}
&
\textbf{Reference Value}
&
\textbf{Proposed Target / Use}
\\
\midrule

Point-cloud representation
&
DP3 \cite{ze2024}
&
512 / 1024 points
&
\textbf{512 points}
\\

Demonstration budget
&
DP3 \cite{ze2024}
&
10--20/task
&
\textbf{10--20 demonstrations/task}
\\

Simulation success
&
Proposed protocol
&
Approximately 92\%
&
\textbf{Approximately 92\% simulation reference}
\\

Real static success
&
Proposed protocol
&
---
&
\textbf{$\geq70\%$ target}
\\

DP3 real success
&
DP3 \cite{ze2024}
&
85\% / 40 demos
&
\textbf{Single-objective external reference only}
\\

Point-cloud advantage
&
DP3 \cite{ze2024}
&
24.2\% relative gain
&
\textbf{Use point-cloud encoder; not a proposed gain target}
\\

Policy/control frequency
&
Diffusion Policy \cite{chi2023}
&
125 Hz
&
\textbf{125 Hz}
\\

Cartesian velocity
&
Diffusion Policy \cite{chi2023}
&
$<0.43$ m/s
&
\textbf{$<0.43$ m/s}
\\

Preference settings
&
Proposed protocol
&
---
&
\textbf{3--5 $\lambda$ settings}
\\

Preference adaptation
&
MORL references
&
No retraining
&
\textbf{No retraining}
\\

Preference objectives
&
Proposed formulation
&
---
&
\textbf{Success, energy, smoothness}
\\

Pareto quality
&
PGMORL / Envelope
&
Baseline HV
&
\textbf{HV $\geq$ PGMORL / Envelope on same wrapper}
\\

GCR-PPO improvement
&
GCR-PPO \cite{munn2025}
&
$\sim$4.55\%; up to $\sim$9.5\%
&
\textbf{Algorithmic order-of-magnitude reference only}
\\

Static safety violation
&
Proposed CMDP
&
---
&
\textbf{Approaches 0\%}
\\

Dynamic safety violation
&
Collision-as-objective
&
Baseline measurement
&
\textbf{Strictly lower than baseline}
\\

\bottomrule
\end{tabular}
\end{table*}

% ============================================================
% 4.4 EXPECTED HYPERVOLUME COMPARISON
% ============================================================

\subsection{Expected Hypervolume Comparison}

Hypervolume is the primary metric for evaluating the quality of the
learned Pareto front.

The comparison will use the \textbf{same vector-reward UR5 benchmark
wrapper} for all principal MORL methods.

\begin{table*}[t]
\centering
\caption{Expected hypervolume comparison under the same visual
vector-reward benchmark.}
\label{tab:expected_hv}
\renewcommand{\arraystretch}{1.3}
\setlength{\tabcolsep}{5pt}

\begin{tabular}{l c c c c l}
\toprule
\textbf{Method}
&
\textbf{Vision}
&
\textbf{Number of Policies}
&
\textbf{Preference}
&
\textbf{CMDP}
&
\textbf{Expected Role}
\\
\midrule

Visual Scalarized PPO
&
\checkmark
&
Multiple
&
Fixed
&
No
&
Visual scalarized baseline
\\

PGMORL
&
\checkmark
&
Multiple
&
Yes
&
No
&
MORL Pareto reference
\\

Envelope MOQ
&
\checkmark
&
1
&
Yes
&
No
&
Single-policy preference reference
\\

\textbf{Proposed Method}
&
\textbf{\checkmark}
&
\textbf{1}
&
\textbf{Continuous $\lambda$}
&
\textbf{Yes}
&
\textbf{Complete framework}
\\

\bottomrule
\end{tabular}
\end{table*}

The expected relationship is

\begin{equation}
\boxed{
HV_{\mathrm{Proposed}}
\geq
\max
\left(
HV_{\mathrm{PGMORL}},
HV_{\mathrm{Envelope}}
\right)
}
\end{equation}

when evaluated using the same vector-reward wrapper and benchmark.

The scalarized PPO baseline is explicitly \textbf{visual}. It uses the
same visual observation setting so that the comparison isolates the
effect of multi-objective learning rather than conflating MORL with
the use of vision.

The hypervolume comparison is therefore the main criterion for judging
whether the proposed MORL formulation provides a useful Pareto
advantage.

% ============================================================
% 4.5 EXPECTED REAL-ROBOT RESULTS
% ============================================================

\subsection{Expected Real-Robot Results}

The final experiment transfers the learned policy from simulation to a
real UR5.

\begin{table*}[t]
\centering
\caption{Expected real-UR5 results.}
\label{tab:real_robot}
\renewcommand{\arraystretch}{1.25}

\footnotesize

\begin{tabular}{l l}
\toprule
\textbf{Metric}
&
\textbf{Expected Real-UR5 Result}
\\
\midrule

Control frequency
&
\textbf{125 Hz}
\\

Policy frequency
&
\textbf{10 Hz + interpolation}
\\

Maximum Cartesian velocity
&
\textbf{$<0.43$ m/s}
\\

Preference settings
&
\textbf{3--5 $\lambda$ settings}
\\

Preference change
&
\textbf{No retraining}
\\

Preference objectives
&
\textbf{Success, energy, smoothness}
\\

Real static success
&
\textbf{$\geq70\%$ target}
\\

Static safety violations
&
\textbf{Approximately 0\% target}
\\

Dynamic safety violations
&
\textbf{Lower than collision-as-objective baseline}
\\

Preference-dependent behavior
&
\textbf{Measurably different trade-offs}
\\

Sim-to-real difference
&
\textbf{Success, energy, smoothness, and safety reported separately}
\\

\bottomrule
\end{tabular}
\end{table*}

The DP3 85\% real-success result is retained as an external
single-objective reference. It should \textbf{not} be presented as the
expected success rate of the complete proposed MORL system.

The proposed real-robot target is instead

\begin{equation}
\boxed{
\text{Real static success} \geq 70\%
}
\end{equation}

while the simulation reference is approximately 92\%.

% ============================================================
% 4.6 WHAT WOULD COUNT AS SUCCESS?
% ============================================================

\subsection{What Would Count as Success?}

To make the evaluation easy to understand, the proposed method will be
considered to meet its principal research objectives if the following
conditions are observed.

\begin{table*}[t]
\centering
\caption{Success criteria for the proposed study.}
\label{tab:success_criteria}
\renewcommand{\arraystretch}{1.25}

\begin{tabular}{l l}
\toprule
\textbf{Condition}
&
\textbf{Target}
\\
\midrule

Visual foundation
&
Approximately 92\% simulation success and $\geq70\%$ real static
success
\\

MORL optimization
&
Improved Pareto quality relative to visual scalarized PPO
\\

Preference conditioning
&
3--5 $\lambda$ values produce measurable behavior changes
\\

No retraining
&
Preference changes use the same trained policy
\\

Pareto quality
&
Hypervolume matches or exceeds PGMORL / Envelope on the same benchmark
\\

Safety
&
Static violation rate approaches 0\%
\\

Dynamic safety
&
CMDP has lower violation rate than collision-as-objective
\\

Real robot
&
Preference-dependent success--energy--smoothness trade-offs remain
observable
\\

\bottomrule
\end{tabular}
\end{table*}

These criteria convert the research questions into measurable
experimental tests rather than qualitative claims.

% ============================================================
% 4.7 IMPORTANT INTERPRETATION
% ============================================================

\subsection{Important Interpretation of the Expected Numbers}

The published numbers and the proposed numbers have different meanings.

\begin{table*}[t]
\centering
\caption{Interpretation of numerical values used in the study.}
\label{tab:number_interpretation}
\renewcommand{\arraystretch}{1.25}

\begin{tabular}{l l}
\toprule
\textbf{Type}
&
\textbf{Meaning}
\\
\midrule

\textbf{Reported}
&
Already measured and reported by the cited paper
\\

\textbf{Reference}
&
Used to establish the performance level or design choice of an
existing method
\\

\textbf{Target}
&
Experimental threshold or configuration proposed for the present
study
\\

\textbf{Expected}
&
Hypothesized behavior of the proposed system
\\

\textbf{Measured}
&
To be filled after the actual experiments
\\

\bottomrule
\end{tabular}
\end{table*}

Therefore, values such as \textbf{4.55\%}, \textbf{9.5\%}, \textbf{85\%},
\textbf{24.2\%}, \textbf{125 Hz}, and \textbf{0.43 m/s} must not be
presented as measured results of the proposed method before the
experiments are performed.

In particular:

\begin{itemize}
    \item The \textbf{85\% DP3 real-success result} is a
    single-objective external reference.

    \item The proposed real static success target is
    \textbf{$\geq70\%$}.

    \item The approximately \textbf{92\% simulation success} is the
    simulation reference for the proposed visual manipulation setup.

    \item The \textbf{24.2\% DP3 value} supports the point-cloud
    representation choice and is not a proposed-system gain target.

    \item The \textbf{4.55\% and 9.5\% GCR-PPO values} provide an
    algorithmic order-of-magnitude reference only and are not UR5
    return targets.

    \item The primary MORL comparison is based on
    \textbf{hypervolume and Pareto quality} relative to PGMORL,
    Envelope MOQ, and visual scalarized PPO.

    \item The preference vector contains only
    \textbf{success, energy, and smoothness}.

    \item \textbf{Safety is handled independently through CMDP
    constraints}.

    \item The safety criterion is \textbf{lower violation rate},
    not a claim that violations can always be completely eliminated.

    \item Static safety violations are expected to
    \textbf{approach 0\%}, while dynamic violations are expected to
    be \textbf{strictly lower than the collision-as-objective
    baseline}.
\end{itemize}

After the experiments, the final paper will replace the expected
values with the actual measured values and report the corresponding
improvement or degradation relative to each baseline.
% ============================================================
% REFERENCES
% ============================================================

\begin{thebibliography}{15}

\bibitem{deheuvel2025}
J.~de Heuvel \emph{et al.},
``Demonstration-Enhanced Adaptable Multi-Objective Robot Navigation,''
IROS, 2025, arXiv:2404.04857.

\bibitem{munn2025}
M.~Munn \emph{et al.},
``Scalable Multi-Objective Robot RL through Gradient Conflict
Resolution,''
arXiv:2509.14816, 2025.

\bibitem{chi2023}
C.~Chi \emph{et al.},
``Diffusion Policy: Visuomotor Policy Learning via Action Diffusion,''
RSS, 2023; IJRR, 2024.

\bibitem{ze2024}
Y.~Ze \emph{et al.},
``3D Diffusion Policy: Generalizable Visuomotor Policy Learning via
Simple 3D Representations,''
RSS, 2024.

\bibitem{xu2020}
J.~Xu \emph{et al.},
``Prediction-Guided Multi-Objective Reinforcement Learning for
Continuous Robot Control,''
ICML, 2020.

\bibitem{yang2019}
R.~Yang \emph{et al.},
``A Generalized Algorithm for Multi-Objective Reinforcement Learning
and Policy Adaptation,''
NeurIPS, 2019.

\bibitem{he2023}
X.~He \emph{et al.},
``Personalized Robotic Control via Constrained Multi-Objective
Reinforcement Learning,''
Neurocomputing, 2023.

\bibitem{abdolmaleki2022}
A.~Abdolmaleki \emph{et al.},
``A Distributional View on Multi-Objective Policy Optimization,''
NeurIPS, 2022.

\bibitem{felten2023}
F.~Felten \emph{et al.},
``MO-Gymnasium: A Toolkit for Reliable Benchmarking and Research in
Multi-Objective Reinforcement Learning,''
NeurIPS, 2023.

\bibitem{pushp2025}
D.~Pushp \emph{et al.},
``Navigating the Wild: Pareto-Optimal Visual Decision-Making in Image
Space,''
arXiv:2511.07750, 2025.

\bibitem{mossalam}
H.~Mossalam \emph{et al.},
``Multi-Objective Deep Reinforcement Learning.''

\bibitem{zhao2023}
T.~Zhao \emph{et al.},
``Learning Fine-Grained Bimanual Manipulation with Low-Cost Hardware
(ALOHA / ACT),''
RSS, 2023.

\bibitem{fu2024}
Z.~Fu \emph{et al.},
``Mobile ALOHA: Learning Bimanual Mobile Manipulation with Low-Cost
Whole-Body Teleoperation,''
CoRL, 2024.

\bibitem{chi2024}
C.~Chi \emph{et al.},
``Universal Manipulation Interface (UMI),''
RSS, 2024.

\bibitem{jiang2025}
Y.~Jiang \emph{et al.},
``IRIS: An Immersive Robot Interaction System,''
CoRL, 2025.

\end{thebibliography}

\end{document}